\chapter{CONTEXTUALIZACIÓN}

\section{Planteamiento del problema}

En el contexto de la industria gastronómica colombiana, los sistemas de gestión de comandas\footnote{Según la RAE, «comanda» (del latín \textit{commanda}) se refiere a un \textit{«documento en que se encarga la ejecución de algo»}, específicamente en el contexto restaurantero sería \textit{«la nota donde se apuntan los pedidos de los clientes»} \cite{rae2024}.} siguen siendo predominantemente basados en procedimientos manuales, donde el personal de meseros transmite las órdenes de los clientes mediante notas escritas a mano que son físicamente entregadas al área de cocina. Aunque este método ha sido históricamente aceptado, su naturaleza primitiva genera desafíos operativos significativos. Entre los principales inconvenientes destacan los retrasos en la ejecución de platos debido a malentendidos en la lectura de las anotaciones, la pérdida accidental de comandas durante el traslado y errores en la facturación derivados de omisiones o doble registro de los ítems.

% --- Párrafo de contexto colombiano ---
Este escenario cobra especial relevancia al considerar la magnitud del sector en el país. En Colombia, existen aproximadamente 167,000 establecimientos gastronómicos registrados \cite{acodres_estadisticas_2021}, donde las Pequeñas y Medianas Empresas (PYME) constituyen la gran mayoría. Según cifras de la Asociación Nacional de Empresarios de Colombia (ANDI), el sector de alojamiento y servicios de comida representa el 3.9\% del Producto Interno Bruto (PIB) nacional, con un aporte de 56.7 billones de pesos en 2022 \cite{andi_pib_2024}. La Cámara del Sector Gastronómico de la ANDI, por sí sola, agrupa a más de 2,038 restaurantes que generan más de 34,822 empleos directos \cite{andi_gastronomia_2023}, lo que subraya la importancia de optimizar sus procesos operativos para garantizar su sostenibilidad y crecimiento.

Estas deficiencias operativas se traducen en cifras concretas. Estudios comparativos indican que los sistemas tradicionales basados en papel presentan una tasa de error en los pedidos de hasta un 8\%, en contraste con el 2\% de los sistemas digitales \cite{Daniel2024}. La causa principal de estas discrepancias se atribuye a errores humanos, como la mala interpretación de la caligrafía o el uso de abreviaturas, que son comunes en entornos de ritmo rápido \cite{LSRetail2025}. Según la Asociación Nacional de Restaurantes de Estados Unidos, este tipo de fallos puede generar pérdidas que ascienden a aproximadamente el 4\% de las ventas totales de un restaurante \cite{Khan2023}. El impacto económico es considerable, ya que se estima que el 60\% de los errores resultan en reemplazos gratuitos que el establecimiento debe asumir \cite{Checkmate2025}.

A pesar de las mejoras tecnológicas disponibles, el persistente uso de comandas físicas se atribuye en parte a la barrera económica que representan los sistemas digitales, cuya implementación requiere una inversión inicial elevada en hardware, software y capacitación, sumada a costos de mantenimiento continuos. Esta situación afecta particularmente a pequeños y medianos establecimientos, para los cuales la digitalización puede resultar poco accesible en comparación con métodos tradicionales de bajo costo operativo \cite{marcelo2024desafios}.

% --- Párrafo sobre vulnerabilidad de PYMES ---
Esta barrera es particularmente crítica en un sector con una alta tasa de mortalidad empresarial. Según la Asociación Colombiana de la Industria Gastronómica (ACODRES), se estima que cuatro de cada diez restaurantes en el país cierran durante sus primeros seis meses de operación, una estadística que resalta una vulnerabilidad estructural donde la ineficiencia operativa puede comprometer la viabilidad del negocio \cite{acodres_estadisticas_2021}. La presión por reducir costos fijos y variables lleva a que muchas PYME gastronómicas posterguen inversiones en tecnología, sin considerar que las pérdidas derivadas de errores manuales pueden superar con creces el costo de una digitalización planificada.

Estas deficiencias no solo impactan la eficiencia operativa de los procesos, sino que también inciden negativamente en la experiencia del cliente, al aumentar los tiempos de espera para la recepción de sus órdenes. La capacidad de realizar modificaciones o correcciones en los pedidos de forma rápida y precisa se ve restringida, lo que puede traducirse en pérdidas económicas tanto a corto como a largo plazo. Consecuentemente, esta problemática afecta la fidelización de los clientes y evidencia fallas en el sistema de gestión de comandas, comprometiendo la competitividad y sostenibilidad de los establecimientos. Asimismo, los incrementos en los tiempos de espera repercuten en la satisfacción del cliente y pueden derivar en notorias pérdidas financieras para el negocio.

Paralelamente a los desafíos operativos, el uso persistente de comandas físicas introduce un riesgo sanitario documentado que trasciende el ámbito de la eficiencia y compromete directamente la inocuidad alimentaria. Stephens et al. (2019) revisaron de manera sistemática la literatura científica sobre los fómites en el entorno construido, concluyendo que constituyen un mecanismo de transmisión relevante para numerosos microorganismos, en paralelo al contacto directo y a la transmisión aérea, y que las características del material del soporte físico inciden directamente sobre la supervivencia microbiana y la tasa de transferencia hacia y desde los humanos \cite{Stephens2019}. En el entorno de un restaurante, la comanda de papel recorre un circuito de múltiples manos: mesero, cajero y cocinero, sin que ninguno de estos puntos de transferencia contemple protocolos de desinfección intermedios, convirtiéndola en un eslabón activo dentro de la cadena de contaminación. Kirchner et al. (2021) demostraron precisamente que las superficies atípicas en entornos de servicio de alimentos, es decir, aquellas que habitualmente quedan fuera de los protocolos de limpieza convencionales, actúan como vectores activos de contaminación cruzada, y que la ausencia de rutinas de higiene específicas para este tipo de objetos perpetúa la transferencia de microorganismos entre el personal y las áreas de preparación de alimentos \cite{Kirchner2021}.

El riesgo particular que representa el papel como vehículo de transmisión ha sido comprobado experimentalmente. Hübner et al. (2011) inocularon papel de oficina estándar con cuatro cepas bacterianas de referencia \textit{Escherichia coli}, \textit{Staphylococcus aureus}, \textit{Pseudomonas aeruginosa} y \textit{Enterococcus hirae} y evaluaron su supervivencia durante siete días bajo condiciones ambientales normales. Los resultados demostraron que todos los microorganismos sobrevivieron en la superficie por un mínimo de 72 horas; de hecho, cepas como \textit{P. aeruginosa} y \textit{E. hirae} mantuvieron niveles de contaminación tan altos que, incluso al séptimo día de prueba, el material nunca alcanzó los umbrales mínimos para considerarse desinfectado. Para evaluar el riesgo de transmisión cruzada, la investigación reprodujo el ciclo de contacto mano--papel--mano, logrando recuperar bacterias viables en la segunda mano en la totalidad de los ensayos independientes, lo que confirma que el papel actúa como un intermediario altamente efectivo para la transferencia de patógenos. Además, los autores destacaron que la naturaleza porosa de este material presenta una resistencia inherente a los agentes químicos líquidos tradicionales. A diferencia de superficies no porosas como el acero inoxidable o el plástico, el papel se convierte en un reservorio persistente que no puede ser neutralizado mediante los métodos de higiene de uso habitual en cocinas \cite{Huebner2011}.

Esta vulnerabilidad adquiere una dimensión adicional cuando se considera la diversidad de agentes patógenos susceptibles de transmitirse a través de objetos de manipulación frecuente. Stephens et al. (2019) señalan que las nuevas herramientas de análisis molecular han evidenciado que los fómites también pueden vehiculizar comunidades microbianas completas y bacterias con resistencia a antibióticos, cuya transferencia indirecta en entornos no clínicos representa un riesgo emergente de especial relevancia en salud pública \cite{Stephens2019}. La relevancia de este riesgo se confirma en el contexto colombiano: un análisis microbiológico en restaurantes formales e informales y puestos de comida en Cali evidenció contaminación cruzada con altos porcentajes de bacterias mesófilas aerobias en superficies de contacto con alimentos, identificando géneros de la familia Enterobacteriaceae incluso en establecimientos de categoría formal \cite{caro-hernandez_analisis_2020}. Esta situación contraviene directamente la Resolución 2674 de 2013 del Ministerio de Salud y Protección Social, que exige a los establecimientos minimizar cualquier foco de contaminación que ponga en riesgo la inocuidad de los alimentos \cite{minsalud_resolucion2674_2013}.


%Paralelamente a estos desafíos operativos, emerge una preocupación crítica de carácter sanitario. El uso continuo de medios manuales, como lápiz y papel, para la comunicación entre áreas expone a riesgos de contaminación cruzada. En entornos donde el manejo frecuente de papeles por múltiples actores ocurre sin garantías de higienización intermedia, se facilita la transmisión de agentes patógenos \cite{Kirchner2021}. Esta situación, que ha cobrado especial relevancia en el contexto postpandemia, obliga tanto a consumidores como a reguladores a revaluar los estándares de seguridad e higiene en la cadena alimentaria. La carencia de un sistema digitalizado no solo limita la optimización de procesos, sino que también constituye una vulnerabilidad en términos de salud pública, comprometiendo la reputación del establecimiento y su cumplimiento con las normativas vigentes.


%La relevancia de este riesgo se confirma con estudios realizados en el contexto colombiano. Un análisis microbiológico en restaurantes formales e informales y puestos de comida en Cali, Colombia, evidenció "contaminación cruzada, con altos porcentajes en bacterias mesófilas aerobias" en superficies de contacto con alimentos \cite{caro-hernandez_analisis_2020}. Aunque la mayoría de los recuentos de coliformes totales se encontraron dentro del límite permisible, la identificación bacteriana demostró la presencia de varios géneros de la familia Enterobacteriaceae. Esta situación contraviene directamente las directrices establecidas en la Resolución 2674 de 2013 del Ministerio de Salud y Protección Social, que regula las buenas prácticas de manufactura y exige a los establecimientos minimizar cualquier foco de contaminación que ponga en riesgo la inocuidad de los alimentos \cite{minsalud_resolucion2674_2013}.


Adicionalmente a los desafíos operativos y sanitarios, el uso de sistemas manuales restringe la capacidad para gestionar cambios dinámicos en los pedidos y llevar un registro estructurado de la información generada. La imposibilidad de actualizar órdenes en tiempo real (sea para ajustar cantidades, modificar ítems del menú, incorporar preparaciones específicas o considerar alergias) propicia errores y confusiones que prolongan los tiempos de atención, repercutiendo negativamente en la satisfacción del cliente \cite{Ser2024Estrategia}. Igualmente, la falta de un registro sistematizado de los alimentos consumidos impide un análisis adecuado de tendencias de demanda, lo cual afecta la optimización de inventarios y la reducción de desperdicios, elementos esenciales para la sostenibilidad económica y ambiental de los establecimientos.

Por otro lado, es preciso destacar que uno de los limitantes más significativos que enfrentan algunos restaurantes al considerar la implementación de software de control de comandas radica en la incertidumbre respecto a la privacidad y seguridad de sus datos. Las soluciones comerciales predominantes dependen, en gran medida, del alojamiento de sus servidores en infraestructuras gestionadas por la misma empresa proveedora del servicio. Este modelo de servicio no contempla, de forma explícita, medidas robustas de encriptación ni garantías sobre el manejo seguro de la información sensible del restaurante. La carencia de certificaciones o protocolos de seguridad específicos genera desconfianza en el operador, ya que el control y la protección de datos vitales como historiales de ventas, inventarios y preferencias de los clientes, queda en manos de terceros. Esta situación representa no solo un riesgo de vulneración de información, sino que además limita la adopción de tecnologías digitales por parte de pequeños y medianos establecimientos, que perciben en este aspecto una barrera adicional frente a la modernización de sus procesos.

En el contexto de la digitalización, esta pregunta problemática invita a analizar de manera precisa cómo el uso de sistemas manuales puede limitar la eficiencia de los procesos operativos, comprometer la seguridad sanitaria y exponer la información sensible en restaurantes de pequeña y mediana escala. Se busca establecer una relación directa entre las deficiencias inherentes a los métodos tradicionales y los inconvenientes que estos generan, para así orientar el desarrollo de una solución tecnológica accesible y efectiva.



\section{Pregunta problema}

¿Cómo desarrollar una aplicación web de bajo costo para la gestión de comandas en restaurantes que garantice la seguridad y confidencialidad de la información mediante encriptación y llaves físicas de acceso?

\section{Objetivos}

\subsection{Objetivo General}
Desarrollar una aplicación web de creación y gestión de comandas para restaurantes que garantice la seguridad y confidencialidad de la información mediante mecanismos de encriptación y el uso de llaves físicas para el acceso administrativo, manteniendo un bajo costo de implementación y alta disponibilidad.

\subsection{Objetivos Específicos}
\begin{itemize}
	\item Diseñar la arquitectura general del sistema, estableciendo los requerimientos funcionales y no funcionales, así como los lineamientos de seguridad y escalabilidad.
	
	\item Implementar la solución tecnológica mediante la codificación y configuración de sus módulos principales, integrando el sistema de autenticación de doble factor basado en NFC con \textit{Raspberry Pi} y módulo \textit{RFID-RC522}.
	
	\item Evaluar el desempeño funcional y la seguridad del sistema mediante pruebas de laboratorio que midan tiempos de respuesta, tasas de éxito en operaciones de comandas, tiempo de autenticación NFC, y efectividad de los mecanismos de encriptación y control de acceso, verificando el cumplimiento de los requerimientos establecidos.
	
	\item Generar la documentación técnica y de usuario del sistema, comprendiendo manual técnico de instalación y configuración, y manual de usuario por roles, facilitando la transferencia de conocimiento y sostenibilidad del proyecto.
\end{itemize}



\section{Justificación}

La persistencia de métodos manuales en la gestión de comandas representa una oportunidad técnica para optimizar la comunicación interna y la experiencia del cliente mediante digitalización. Este proyecto desarrolla una solución web de bajo costo que transforma el proceso de gestión de pedidos en restaurantes PYME, haciéndolo accesible para establecimientos que enfrentan limitaciones económicas. Al eliminar la barrera de inversión inicial elevada característica de sistemas comerciales, la propuesta ofrece escalabilidad y rentabilidad sin comprometer funcionalidad, convirtiéndose en una alternativa viable para negocios de distintos tamaños que buscan modernizar sus operaciones con mínimo impacto financiero.

El contexto colombiano presenta una oportunidad significativa para esta solución. Según el DANE, la Encuesta de Micronegocios (EMICRON) del primer trimestre de 2024 evidencia la persistencia de micronegocios como motor económico del país \cite{DANE2024}, sector que enfrenta limitaciones estructurales de inversión tecnológica. Los beneficios técnicos y operativos del sistema se estructuran en seis ejes fundamentales:

\begin{itemize}
	\item \textbf{Estandarización del Flujo de Información:}   
	El sistema establece un canal digital unificado que conecta sala, cocina y caja mediante registros estructurados, eliminando la variabilidad inherente a las notas manuscritas. Investigaciones recientes demuestran que los sistemas tradicionales de pedidos son labor-intensivos, resultando en altas tasas de inexactitud en órdenes \cite{Sultana2024}, mientras que sistemas automatizados minimizan estos errores al asegurar colocación precisa y exacta de pedidos \cite{Sultana2024}. Al agilizar la comunicación entre clientes y personal de cocina, el sistema reduce errores, minimiza tiempos de procesamiento de órdenes y mejora la eficiencia del flujo de trabajo \cite{Katagri2025}.
	
	En el estudio de caso \textit{Using Service Blueprinting to Analyze Restaurant Service Efficiency} de Emily Hummel y Kevin S. Murphy\cite{EmilyHummel}, se demuestra que la eficiencia operativa mejora cuando la información se comunica de manera efectiva y estandarizada. El sistema propuesto logra precisamente esto al reemplazar la comunicación verbal y escrita, que es propensa a errores, por un registro digital único y consistente.
	
	
	\item \textbf{Trazabilidad Operativa en Tiempo Real:} 
	Más allá de la simple transmisión de pedidos, el sistema implementa capacidades de monitoreo temporal que registran cada modificación con marca de tiempo, identifica al responsable de cada operación y calcula métricas de servicio desde la toma del pedido hasta su entrega. Estudios académicos confirman que la automatización de la gestión de pedidos reduce el tiempo de procesamiento en aproximadamente 40\%, mientras que la eliminación de inexactitudes debidas a errores manuales mejora significativamente la satisfacción del cliente basada en retroalimentación de usabilidad \cite{Katagri2025}.
	
	Según De Vries, Roy y De Koste en el artículo \textit{Worth the wait? How restaurant waiting time influences customer behavior and revenue} \cite{DEVRIES201859}, la capacidad de los sistemas digitales para agilizar procesos de pedido y pago impacta directamente en la retención de clientes, la frecuencia de visitas y la duración de estancias \cite{DEVRIES201859}. El control temporal preciso que ofrece esta solución convierte estos beneficios en métricas medibles y gestionables.
	
	\item \textbf{Mitigación de Riesgos Sanitarios:} 
	La eliminación del papel como medio de transmisión de información entre áreas reduce vectores potenciales de contaminación cruzada. Si bien la literatura sobre contaminación específica en comandas es limitada, estudios sobre menús físicos demuestran que superficies manipuladas repetidamente sin protocolos de higienización sistemática presentan cargas microbiológicas significativas. 
	
	Investigaciones han evidenciado que los menús en papel, al ser manipulados repetidamente por clientes y personal, presentan altos niveles de contaminación \cite{alsallaiy2015}, y que la evaluación visual resulta insuficiente para garantizar limpieza \cite{choi2025}. Dado que las comandas manuscritas atraviesan un flujo similar de manipulación múltiple, la digitalización representa una medida preventiva al eliminar completamente este contacto físico entre áreas operativas.
	
	\item \textbf{Arquitectura de Seguridad Multicapa:} 
	La protección de información del establecimiento se implementa mediante dos mecanismos complementarios: encriptación de base de datos y autenticación física para acceso administrativo. La encriptación resguarda datos sensibles frente a accesos no autorizados incluso en caso de compromiso del servidor, mientras que la llave física desarrollada específicamente para este proyecto establece una barrera de acceso tangible que complementa controles digitales. Esta combinación de seguridad lógica y física genera un nivel de protección integral adaptado a las crecientes amenazas en entornos digitales, factor determinante para la confianza y adopción tecnológica en contextos donde la información operativa y financiera requiere máxima confidencialidad.
	
	\item \textbf{Integración Interdisciplinaria Hardware-Software:} 
	El proyecto constituye una aplicación integral de Ingeniería Electrónica al desarrollar un sistema embebido de autenticación mediante tecnología NFC acoplado a la aplicación web. La implementación de esta llave física de autenticación requiere conocimientos específicos de sistemas embebidos, protocolos de comunicación inalámbrica de corto alcance a frecuencia de 13.56 MHz conforme al estándar ISO/IEC 14443 \cite{Liyanage2018}, y arquitecturas de seguridad multicapa que combinan componentes electrónicos con algoritmos criptográficos. 
	
	Esta convergencia entre hardware especializado y desarrollo de software representa el núcleo de la formación en Ingeniería Electrónica, donde la comprensión de principios electromagnéticos, comunicaciones digitales y sistemas de control se aplican para resolver problemas reales de automatización y seguridad. El diseño del dispositivo NFC como mecanismo de autenticación física complementa controles lógicos implementados en software, demostrando la transversalidad característica de proyectos electrónicos modernos que trascienden el desarrollo puramente informático al integrar sensores, actuadores y sistemas de comunicación de bajo nivel con capas de aplicación de alto nivel \cite{Liyanage2018}.
	
	\item \textbf{Inteligencia Operativa Basada en Datos:} 
	El sistema transforma registros transaccionales en capacidades analíticas mediante el seguimiento sistemático de ingresos, gastos y patrones de consumo. Esta funcionalidad permite identificar tendencias de demanda, optimizar niveles de inventario y sincronizar tiempos de abastecimiento con necesidades reales, reduciendo desperdicio de alimentos con impacto económico y ambiental directo. Adicionalmente, el registro automático de horarios de entrada y salida del personal facilita la gestión de asistencia con datos precisos y confiables.
	
	Basado en la conclusión de Nabeel Salih Ali, Hasanein D. Rjeib, et al., en \textit{Automated attendance management systems: systematic literature review}, los sistemas de gestión de asistencia \textit{"proporcionan datos precisos y confiables, mejoran la eficiencia, facilitan el cumplimiento normativo y ofrecen datos útiles para la toma de decisiones"} \cite{Automatedattendance}, complementando las capacidades de gestión de comandas con control integral de recursos humanos.
\end{itemize}

La arquitectura técnica del sistema prioriza la optimización de recursos, permitiendo su operación en equipos con especificaciones computacionales básicas tanto para servidor como para base de datos. Esta característica elimina la necesidad de inversión en hardware especializado o servicios costosos de hosting en la nube, reduciendo significativamente barreras económicas para pequeños y medianos negocios. Al implementarse como aplicación web, el sistema garantiza compatibilidad universal con dispositivos heterogéneos \cite{BALSAM2025112186}, desde smartphones para toma de pedidos hasta computadoras para supervisión administrativa, independientemente del sistema operativo. La integración mediante red de intranet local elimina dependencias de conectividad externa a internet, suprimiendo gastos operativos recurrentes y reduciendo tanto el presupuesto de implementación inicial como los costos mensuales de operación \cite{akyildiz2006survey}.

En conjunto, estas características posicionan la solución como una herramienta de transformación integral para restaurantes que combina eficiencia operativa, adaptabilidad tecnológica y sostenibilidad económica. La convergencia de capacidades de gestión de pedidos, control de personal, análisis de datos operativos y seguridad multicapa diferencia esta propuesta de sistemas comerciales fragmentados que abordan aspectos aislados. Los indicadores cuantificables de impacto documentados en literatura académica incluyen reducción aproximada del 40\% en tiempo de procesamiento de órdenes y eliminación de inexactitudes por errores manauales \cite{Katagri2025}, métricas que justifican la inversión para el sector PYME colombiano. La incorporación de mecanismos de seguridad tanto sanitarios como informáticos provee garantías sobre confidencialidad y manejo seguro de información, factor decisivo para facilitar adopción tecnológica en un sector tradicionalmente resistente al cambio. El componente de ingeniería electrónica representado por el sistema embebido NFC conforme a estándares ISO/IEC 14443 \cite{Liyanage2018} demuestra la aplicación interdisciplinaria de conocimientos en comunicaciones inalámbricas, sistemas digitales y arquitecturas de seguridad, diferenciando este proyecto de desarrollos puramente informáticos. Estas ventajas consolidan la propuesta como una solución innovadora que impulsa la transformación digital en el sector gastronómico, contribuyendo al crecimiento económico y la sostenibilidad de establecimientos PYME.




